This survey has reviewed the state of the art in multimodal Gaussian
Splatting SLAM along three depth-source axes: LiDAR+RGB+IMU,
iToF+RGB+IMU, and monocular depth+RGB+IMU.  The key findings are:

\begin{enumerate}
    \item \textbf{LiDAR-based systems} (Gaussian-LIC2, RMGS-SLAM, GS-LIVO,
    etc.) achieve the highest accuracy and outdoor robustness through
    direct metric geometry, but at high hardware cost.  The field is
    converging on depth completion networks to address LiDAR sparsity and
    bidirectional map--odometry coupling for improved tracking.

    \item \textbf{iToF-based systems} (ToF-Splatting) are in their infancy
    but demonstrate that even $8\times8$ sparse ToF depth, fused with
    multi-view geometry and monocular cues, can achieve competitive SLAM
    performance---a critical enabler for mobile and AR/VR deployment.
    The integration of IMU into iToF Gaussian SLAM remains an open
    opportunity.

    \item \textbf{Monocular-depth-based systems} (VIGS-SLAM, HI-SLAM2,
    VGGT-GS SLAM, etc.) are the most accessible and rapidly evolving,
    leveraging foundation model depth priors and tight visual-inertial
    coupling to approach RGB-D quality without active depth sensors.
    Scale ambiguity remains the core challenge, addressed by IMU fusion,
    metric depth models, or grid-based scale alignment.
\end{enumerate}

The unifying trend across all three categories is the shift from
\emph{passive mapping} (Gaussians as visualization) to \emph{active
mapping} (Gaussians as geometric constraints for tracking, depth
completion, and loop closure).  As 3DGS representations become lighter
and foundation model priors become more accurate, we expect multimodal
Gaussian Splatting SLAM to transition from research prototype to
practical deployment on consumer devices.
